\section{What DUALEXIS Is NOT}
\label{sec:what-dualxis-is-not}

This page clarifies common concerns raised by schools, families, ethics committees,
data protection officers, and regulators when discussing monitoring in educational settings.
It describes \textbf{design boundaries} for governed pilots---not marketing claims.

\subsection{Explicit Exclusions}

DUALEXIS is \textbf{not}:

\begin{itemize}[noitemsep]
  \item \textbf{Facial recognition} --- it does not identify who someone is from face images.
  \item \textbf{Biometric surveillance} --- it does not use voiceprints, gait fingerprints,
        or other biometric templates.
  \item \textbf{Predictive policing} --- it does not forecast crime, intent, or individual
        future misconduct.
  \item \textbf{Behavioral profiling} --- it does not build psychological, disciplinary,
        or longitudinal profiles of students or staff.
  \item \textbf{Autonomous lockdown enforcement} --- it cannot trigger lockdowns, alarms,
        or physical security actuation on its own.
  \item \textbf{Automated disciplinary scoring} --- it does not assign punishments, grades,
        or conduct scores.
  \item \textbf{Cloud-based raw media surveillance} --- video and audio are not uploaded
        to commercial cloud platforms for routine storage or review in the reference pilot design.
  \item \textbf{Student tracking infrastructure} --- it does not follow named individuals
        across zones, days, or campuses.
  \item \textbf{Identity analytics} --- it does not link events to student IDs, staff directories,
        or national identifiers in operational outputs.
\end{itemize}

These exclusions are \textbf{architectural design invariants} in the reference specification.
They are not optional settings that a school may enable without a separate governance process,
DPIA update, and ethics review.

\subsection{What the Pilot Actually Provides}

Under governed deployment, DUALEXIS offers \textbf{zone-level situational summaries}
(for example: elevated crowding near an exit, acoustic stress in a cafeteria) and
\textbf{advisory text} for trained staff to read, question, and act upon---or dismiss.

That is a different category of system from identity-centric surveillance or automated
enforcement.

\subsection{Authority, Advisory Role, and Institutional Responsibility}

\textbf{Human authority remains central.}
Licensed safety personnel and school leadership decide what each summary means and whether
any protocol step is warranted.

The system is \textbf{advisory only}.
It may suggest that staff \emph{monitor}, \emph{verify}, or \emph{consider escalation};
it does not execute those steps.

\textbf{Operational decisions remain institutional responsibilities.}
Public-address announcements, radio calls, contact with emergency services, parent
notification, and disciplinary processes stay within existing school policies and law.

\textbf{Privacy constraints are built into the design}, not added as an afterthought:
short-lived sensing, no default media archive, no biometrics, structured records suitable
for audit, and processing on school-controlled infrastructure subject to the site's DPIA
and steering committee.

A pilot should proceed only where these boundaries are understood, documented, and
accepted---not where stakeholders expect identification, automation, or cloud video review
by default.
